\section{Introduction}

\subsection{Background}
\setresponsible{Surya Salim}
Large amounts of organic waste are produced every day by the meat-processing industry. This waste is harmful to soil and water if it is not properly treated and disposed of. Anaerobic digestion is one established route for handling this material, where micro-organisms break the material down in the absence of oxygen and release biogas. A byproduct of this process is digestate, which contains most of the nitrogen, phosphorus, and other nutrients from the original feedstock and is typically applied as fertilizer to farmland. However, the amount that can be used in this manner is limited due to its high water content, changeable composition, and expensive transportation and storage costs \parencite{tambone-2010}. The residue of a recycling process, therefore, becomes a waste problem of its own.

\vspace{0.3em}

Pyrolysis turns that residue into a more stable and useful solid. Heating the dried digestate in an oxygen-limited atmosphere drives off the volatile fraction and leaves a porous, carbon-rich material called biochar \parencite{Qi2024}. Biochar can be returned to the soil, where it improves structure, holds water and nutrients, and stores carbon for long periods \parencite{Lehmann2015}; with further treatment, it can also act as a porous adsorbent for removing contaminants from gasses and liquids. It is combustible as well and thus has some value as a solid fuel; however, that is a secondary consideration in the present work.

\vspace{0.3em}

Most of what is known about biochar comes from wood and crop feedstock. Digestate-derived biochar, and slaughterhouse digestate in particular, has received far less attention, even though its mineral-rich composition makes it behave quite differently from lignocellulosic char \parencite{opatokun-2015}. A further practical point is that the digestate studied here arrived without any supplier specification of its composition, which is why its elemental make-up is determined in-house rather than taken from a data sheet.

\vspace{0.3em}

The slaughterhouse case closes a loop. Anaerobic digestion of slaughterhouse waste produces biogas contaminated with hydrogen sulfide (H$_2$S), a toxic and corrosive gas that must be removed before the biogas can be used \parencite{salminen-2002}. Activated biochar made from the digestate can adsorb the H$_2$S. The result is a sulfur-enriched biochar that can then be applied to soil, where sulfur serves as a plant nutrient \parencite{Zhang2017SulfurEnrichedBiochar}. A single feedstock thus addresses two problems at once: the digestate is diverted from disposal, and the H$_2$S is stripped from the biogas, with both ending up in a material that returns nutrients to the field.

\subsection{preceding Project Work}
\setresponsible{Regen Wijaya}
The present work builds on a preceding project in the master's program renewable energy as part of a research project on “Hybrid Biochemical and Thermochemical Conversion of Slaughterhouse Biowaste for Renewable Energy Production”. The digestates were provided by the institutes as part of a collaborative partnership within the research project.

\vspace{0.3em}

The preceding work investigated the production of biochar from three different types of digestates:

\begin{itemize}
    \item \text{Biowaste Fermentation} 
    \item \text{Grass silage}
    \item \text{Slaughterhouse}
\end{itemize}

The resulting biochars were subsequently activated with carbon dioxide, followed by  scanning electron microscopy (SEM) analysis to visualize the morphology and pore development as well as CHNS-analysis to determine the elemental composition of the biochars \parencite{MartinezHuete2024Projektarbeit}. 

\subsection{Motivation of Work}

This work follows the same structure as the preceding one. However, due to several changes and additions, its scope differs slightly. This work focused exclusively on investigating the production and characterization of biochar using slaughterhouse waste originating from France and Algeria, which—as in the preceding project—was provided by the same corporate partner as part of a research project. The French digestate was designated as the primary substrate, since the Algerian digestate was the same digestate that had already been used in the preceding project work. This work examined exclusively the activation of biochar using steam in order to compare the results of steam activation and CO$_2$ activation for biochar derived from slaughterhouse waste. 

\vspace{0.3em}

To determine more properties of the biochar, several characterization methods were employed that had not been used in the preceding study. In this study, the calorific value of the biochar was determined to obtain realistic information about the energy content of the digested residue and the resulting biochar. In addition, the characterization included a Brunauer-Emmett-Teller (BET) analysis to determine the surface area and pore structure of the biochar, as well as a CHNS analysis to determine the elemental composition of the activated biochar.

\subsection{Objectives of The Work}

\subsubsection{General Objective}

The general objective of this work is to explore the potential of different digestates from slaughterhouses for biochar production, including their potential modification through activation and sulfur loading. This work investigates the effect of pyrolysis temperature on biochar production, the effect of steam as an activating agent on biochar activation, and the feasibility of sulfur enrichment on the resulting biochars. This work also aims to characterize the properties of the resulting biochars. In addition, another objective of this work is to draw comparisons with the authors’ preceding project on the production and characterization of biochar via pyrolysis and CO$_2$ activation and to discuss the findings.

\subsubsection{Specific Objectives}

To help achieve the overall objectives of this work, several specific objectives are formulated. These are as follows:

\begin{itemize}
    \item to produce biochar from two different slaughterhouse digestates using pyrolysis with temperatures at 400, 500, and 600~$^\circ$C, and further determine the yield obtained at each temperature. The higher heating value of the char is also determined by bomb calorimetry as a complementary fuel property 
    \item to perform the activation of biochars with steam and to quantify the effect of activation through the associated mass loss.
    \item to carry out the sulfur-loading process on the activated biochars with hhydrogen sulfide.
    \item to characterize the physicochemical properties of the produced biochar using various methods, including the determination of calorific value, BET surface area, and CHNS elemental analysis, and to evaluate potential applications based on the determined properties.  
    \item to compare the result from the current work, and to set the findings against the CO$_2$-activation baseline of the authors’ preceding project work and discuss the finding.
\end{itemize}

\subsubsection{Limitations of the Work}
\label{subsec:limitations}
In this work, several limitations were encountered, particularly in the pyrolysis and steam activation experiments. Firstly, the amount of produced biochars is limited due to the limited amount of biomass/feedstock available to undergo pyrolysis. This could affect further experiments that require a large amount of biochars. A further limitation concerned the steam activation system, in particular the reactor tube and the vessel, which were not in perfect condition. Furthermore, it was not possible to install a suitable water pipe connected directly to the reactor, as this would have been too time-consuming. Therefore, water was introduced into the reactor using a single-use syringe. The syringe injection restricted control over the steam supply. Although the water was injected at a temperature of over 100~$^\circ$C, partial condensation in cooler parts of the reactor may have lowered the steam quantity for activation.

\vspace{0.8em}

A limitation was also encountered in analyzing or characterizing the resulting biochars. Analysis or characterization methods such as BET surface area analysis and CHNS elemental analysis had to be outsourced and were relatively expensive. Therefore, only several biochars were selected. SEM analysis was also not possible due to resource constraints.
\subsection{Structure of The Thesis}

The thesis is structured into several chapters, each of which covers specific aspect of the research that together outline the overall framework of this work as follows:

\vspace{0.8em}

\textbf{chapter~1} introduces the context of the work, consisting of the research background, objectives, and structure of the thesis.

\vspace{0.8em}

\textbf{chapter~2} presents the theoretical background of the work, serving as the basic principles for the experiment and analysis. 

\vspace{0.8em}

\textbf{chapter~3} outlines the experimental investigations that were carried out within the scope of this work, covering the materials as well as the methods for the experiment.

\vspace{0.8em}

\textbf{chapter~4} presents the results of the work. The results are interpreted to highlight the key findings related to the overall objectives of this work. The key findings will be concluded at thhe end.

\vspace{0.8em}

\textbf{chapter~5} presents the conclusion of the experiments.